\begin{abstract}
Ce rapport rassemble, en un seul document, l'état du projet « jumeau numérique de l'offre de soins libérale » au
15 novembre 2026, pour les équipes qui le reprendront : chercheurs, statisticiens et professionnels de la
planification sanitaire.

\paragraph{Le modèle.} On projette jusqu'en 2050 l'effectif du secteur privé (« libéral » au sens large) de cinq
professions : sages-femmes, chirurgiens-dentistes, pharmaciens, masseurs-kinésithérapeutes et infirmiers. La
projection porte sur chaque territoire, de la commune à la France. Elle simule des politiques publiques nationales
exprimées comme des réglages datés d'hyperparamètres sourcés. Le modèle est un jumeau à état explicite
(profession × territoire × âge × sexe × secteur), composé de modules de flux à contrat vérifiable. L'estimateur
servi est un modèle de cohortes (C1).

\paragraph{Le principal résultat de précision.} Au banc départemental à cinq ans, l'erreur absolue moyenne en
pourcentage de C1 est de 7,76\,\%\mesure, contre 9,94\,\%\mesure pour la persistance et 9,32\,\%\mesure pour la
tendance linéaire. Le budget d'erreur attribue 3,11 points\mesure au seul bruit de comptage, qui est
incompressible. Trois familles de pistes n'ont pas battu C1 : des intervalles de prévision recalibrés (prévision
conforme, réconciliation MinT), un gradient boosting global et un perceptron multicouche.

\paragraph{Ce qui n'est pas acquis.} Plusieurs résultats sont négatifs et sont exposés comme tels :
\begin{itemize}
\item aucun chiffre à l'horizon 2050 n'est citable comme prévision ;
\item aucun grain inférieur au département n'est calibré ;
\item aucun levier n'est validé de bout en bout par une évaluation causale ;
\item pour la moitié des protocoles pré-enregistrés, l'antériorité du gel sur le calcul n'est pas prouvable par
  l'historique du dépôt.
\end{itemize}

\paragraph{Le reste du document.} Il couvre aussi l'interface, la couche de langage naturel et l'assistant, la
conformité (RGPD, AIPD, règlement 2017/745), un inventaire chiffré des questions ouvertes et le programme du
module de besoin de soins. Chaque nombre porte sa provenance : mesuré\mesure, dérivé\derive{} ou
supposé\suppose. Les incohérences relevées entre les fichiers du dépôt sont listées en annexe, avec la valeur
retenue.
\end{abstract}
